% =====================================================================
%  DAY 2 — 2026 Comprehensive Exam
%  Ryan Gibbons
%
%  HOW TO USE
%    Write the Day 2 response in the body below.
%    Build:  latexmk -pdf "Day 2 Gibbons.tex"
%      or:   python3 tools/build_all.py "Day 2*"
%
%  CITATIONS — American Political Science Review style.
%    \cite{key}            -> (Putnam 1993)          parenthetical
%    \citeasnoun{key}      -> Putnam (1993)          narrative / in-text
%    \citeyear{key}        -> 1993                   year alone
%    \cite[45]{key}        -> (Putnam 1993, 45)     with page number
%    \possessivecite{key}  -> Putnam's (1993)        possessive
%    Multiple: \cite{a,b,c}
%    Keys come from ../refs.bib. Nothing is cited yet — that is fine;
%    the References section simply stays empty until you cite something.
% =====================================================================


% ---------------------------------------------------------------------
%  SETTING 1 of 3 — FONT SIZE
%  Change the option on the next line. Valid values: 10pt, 11pt, 12pt.
% ---------------------------------------------------------------------
\documentclass[12pt]{article}


% ---------------------------------------------------------------------
%  SETTING 2 of 3 — MARGINS
%  One value applied to all four sides. Examples: 1in, 1.25in, 2.5cm.
%  For uneven margins, replace the \geometry line further down.
% ---------------------------------------------------------------------
\newcommand{\MarginSize}{1in}

% ---------------------------------------------------------------------
%  SETTING 3 of 3 — LINE SPACING
%  1.0 = single, 1.5 = one-and-a-half, 2.0 = double. Decimals are fine.
% ---------------------------------------------------------------------
\newcommand{\LineSpacing}{2.0}


% ---------------------------------------------------------------------
%  TITLE PAGE FIELDS — edit these four lines per response
% ---------------------------------------------------------------------
\newcommand{\PaperTitle}{Day 2}
\newcommand{\PaperAuthor}{Ryan Gibbons}
\newcommand{\PaperDate}{\today}
\newcommand{\PaperSubtitle}{2026 Comprehensive Examination \\[0.3em] Day Two}


% =====================================================================
%  PACKAGES
% =====================================================================

% --- Page layout, spacing, headers ---
\usepackage[margin=\MarginSize]{geometry}
\usepackage{setspace}
\usepackage{fancyhdr}
\usepackage{titlesec}
\usepackage{titling}

% --- Fonts and text ---
% lmodern is Computer Modern redrawn as a scalable outline font: it looks the
% same as the LaTeX default but, unlike the default, microtype can expand it
% and PDF text search/copy works properly.
\usepackage{lmodern}
\usepackage[T1]{fontenc}
\usepackage[utf8]{inputenc}
\usepackage{microtype}      % subtle spacing improvements
\usepackage{csquotes}       % \enquote{...} for nested quotation marks
\usepackage{soul}           % \hl{} highlight, \ul{} underline
\usepackage{xcolor}
\usepackage{url}

% --- Mathematics ---
\usepackage{amsmath}
\usepackage{amssymb}
\usepackage{amsfonts}
\usepackage{amsthm}
\usepackage{mathtools}      % \coloneqq, \DeclarePairedDelimiter, etc.
\usepackage{bm}             % \bm{} bold math
\usepackage{dsfont}         % \mathds{1} indicator function

% --- Figures ---
\usepackage{graphicx}
\usepackage{float}          % [H] to pin a float exactly here
\usepackage{subcaption}
\usepackage{caption}
\usepackage{wrapfig}
\usepackage{rotating}

% --- Tables ---
\usepackage{booktabs}       % \toprule \midrule \bottomrule
\usepackage{tabularx}
\usepackage{longtable}
\usepackage{multirow}
\usepackage{array}
\usepackage{threeparttable}
\usepackage{siunitx}        % \num{1234.5}, aligned numeric columns

% --- Lists, code, diagrams ---
\usepackage{enumitem}
\usepackage{listings}
\usepackage{tikz}
\usetikzlibrary{shapes.geometric, arrows, arrows.meta, calc, positioning, trees}
\usepackage{pgfplots}
\pgfplotsset{compat=1.18}

% --- Citations: APSR style (harvard.sty + apsr.bst, both in this folder) ---
\usepackage[apsr]{harvard}

% --- Cross-references. Loaded last, as hyperref prefers to be. ---
\usepackage{hyperref}
\hypersetup{
  colorlinks = true,
  linkcolor  = black,      % internal links: black for a printed draft
  citecolor  = black,
  urlcolor   = blue,
  pdftitle   = {\PaperTitle},
  pdfauthor  = {\PaperAuthor},
}


% =====================================================================
%  APPLY SETTINGS
% =====================================================================
\setstretch{\LineSpacing}

% Page numbers bottom-right, no rules, no header text.
\pagestyle{fancy}
\fancyhf{}
\rfoot{\thepage}
\renewcommand{\headrulewidth}{0pt}
\renewcommand{\footrulewidth}{0pt}

% Paragraphs: indented, no extra vertical gap (standard for prose).
% For block paragraphs instead, set parindent to 0pt and parskip to ~1em.
\setlength{\parindent}{0.5in}
\setlength{\parskip}{0pt}


% =====================================================================
%  THEOREM ENVIRONMENTS
% =====================================================================
\theoremstyle{plain}
\newtheorem{theorem}{Theorem}
\newtheorem{prop}{Proposition}
\newtheorem{lemma}{Lemma}
\newtheorem{corollary}{Corollary}

\theoremstyle{definition}
\newtheorem{definition}{Definition}
\newtheorem{assumption}{Assumption}
\newtheorem{example}{Example}

\theoremstyle{remark}
\newtheorem*{remark}{Remark}


% =====================================================================
%  DOCUMENT
% =====================================================================
\begin{document}

% --- Title page (no page number) ---
\begin{titlepage}
    \centering
    \vspace*{2cm}
    {\large 2026 Comparative Politics Comprehensive Exam} \\
    \vspace{1cm}
    {\large Day 2: Comparative Institutions} \\
    \vspace{8cm}
    {\large \PaperAuthor \par}
    \vspace{.2cm}
    {\large University of Chicago \par}
    \vspace{.2cm}
    {\large Department of Political Science \par}
    \vspace{1cm}
    {\large \PaperDate \par}
    \vfill
\end{titlepage}


\newpage
\thispagestyle{empty}

\vspace{2cm}

\section*{Question}

Are institutional constraints on action always necessary for the resolution of social/political/economic problems or for the stability of social/political/economic relations? Consider various reasons given in the literature for the need for institutions. How do institutions come into being? Is it possible to have social cooperation in politics without institutions? Provide empirical examples for different approaches to the emergence and design of institutions

% Body starts here; page numbering restarts at 1.
\newpage
\setcounter{page}{1}

\section*{Introduction}



In this response, I examine the primary factors determining whether or not social, political, and economic phenomena are most appropriately addressed through institutional constraints or non-institutional alternatives, and how these ``solutions'' come to exist. To do so, we must first construct a justification \emph{for} institutions as a means of securing cooperation, and then identify the cases under which other solutions may be viable. The paper proceeds as follows. First, I discuss the primary purposes of institutions in facilitating cooperation, examining the general frameworks which they may be used to solve and how they arise from those frameworks. Second, I explore how repetition of interactions may provide non-institutional stability under certain conditions (which themselves may be secured by institutions). Third, I consider how existing incentive structures may already be aligned with optimal outcomes, transitioning into a focus on the role of social norms as informal institutions. I conclude by discussing the ever-present philosophical debate between rational-choice-focused ``new institutionalists'' and their interpretivism-focused dissenters, and consider how our framing of institutions might facilitate a reconciliatory (and intellectually productive) future between these two factions.

I am, both in head and heart, an institutionalist. Yet, I concede that the degree to which institutions are necessary to facilitate the effective operation of society is not uniform. In this response, I intend to characterize a relatively broad literature within institutional theory, and consider the limits of its applicability. If my experience with the preceeding section of this examination has taught me anything, it is that I cannot reasonably hope to discuss every such factor in the coming hours. Instead, I must relegate myself to exploring them in the coming years and decades.


\section*{Institutions and Cooperation}


North's \citeyear{north1990} construction of institutions as ``the rules of the game'' is perhaps one of the most famous definitions in contemporary political science, and for good reason. We can broadly characterize institutions as those influences which shape outcome behavior in a given social setting. We can map two such means of influence from North's charcterization: either the rules are unbreakable, in which they ultimately define the action set available, or they are breakable but utility-changing, in which they affect the \emph{utilities} associated with the available action set (and thereby, under a rationalistic framework, the outcome behaviors selected). In either case, the end effect -- behavioral alteration -- can be achieved, though we will explore the limitations of this definition in our discussion of existing incentive alignment.

Yet, in determining when institutions are necessary, we must examine not only what they \emph{are}, but also what they \emph{do}.\footnote{or ought to do} Let us first consider North's \citeyear{north1990} characterization of institutions as ``reducing transaction costs.'' Through this framework, institutions exist to counter market\footnote{This is not merely an economic or trade market, but very well may include a political or policy market} inefficiencies and facilitate lower-cost exchanges of goods, information, and power. There are a variety of ways in which institutions can deliver this efficiency, the first being as an informational structure: institutions provide stability in expectations such that actors can reliably predict the consequences of their behavioral decisions \cite{mahoneythelen2009}. This reduction in uncertainty can prove a valuable asset to efficient market interactions, and can avoid the losses that come with calculated caution \cite{handleylimao2015}.

Though these above mechanisms invoke clear reductions in transactional costs, at least insofar as ``traditional'' market exchanges, this conceptualization can be expanded to three further core problems in political science. Each could easily warrant multiple exams' worth of examination; I instead attempt to provide a quick-yet-definitive characterization of each and the institutional responses therein.

\subsection*{Institutions for Coordination}

A coordination problem involves a scenario whereby actors have multiple equilibria which they can collectively achieve, but face obstacles to deciding which equilibrium\footnote{encompassing both outcome and process} to target. The canonical example under symmetry is a society's decision to drive on either side of the road. Presumably, most preferences for one side over the other are relatively weak, but it is clear that all rational members of the society should prefer a uniform decision: either everybody drives on the left, or everybody drives on the right, lest we risk crashes and inefficiency. Here, a reduction in transaction costs (car crashes) can be supported by institutions which either (1) facilitate communication such that an agreed-upon outcome can be decided, or (2) outright impose one of the two equilibria (``all left'' or ``all right''). Legislative institutions often serve the first purpose \cite{gandhiprzeworski2007}, while the second is particularly ripe for political entrepreneuriship \cite{taylor1990}. However, each of these functions still requires an initial instigation; a legislature does not exist until it does, and a political entrepreneur must actually accumulate support as a coordinative device.

Yet, coordination problems need not feature this symmetry, and indeed it is quite frequent within the political sphere for actors to prefer a coordinated outcome, but to diverge in their preferences \emph{between} coordinated outcomes. A common example is within ideological factions, whereby the meaningful representation of minority groups may require their cooperation, but the terms of that cooperation are not necessarily given. Consider a three-party environment where no party holds an outright majority (as if one does, the question is largely moot). By definition, any two parties could form a coalition to gain office, but in doing so would almost certainly need to compromise on some policy dimensions \cite{strom1990,mullerstrom1999,cheibubetal2004}. This tradeoff, as Boix \citeyear{boix1999} demonstrates, can produce incredibly counterproductive results if a clear coordination mechanism is not present; consolidation into a majority coalition can be nearly impossible if the relevant parties are evenly strong such that one of the two coordinating equilibria can be clearly selected. Accordingly, the party organizations themselves exist as institutions to solve one coordination problem, but when considering ideological or policy coalitions, further institutional involvement may be warranted.

How, then, do these coordinative institutions emerge? Some scholars emphasize the role of exogenous and spontaneous formation of such coordinative institutions which, upon ``solving'' the previous coordination problem, are retained \cite{sugden1989,kuran1991}. Yet, this suggestion of exogeneity\footnote{As suggested in yesterday's response -- not a particularly appealing explanation in social science} was quickly met with literature focusing on their endogenous formation, even if fueled in part by certain circumstances. Aldrich's \citeyear{aldrich1995} \emph{Why Parties} provided a highly-intentional narrative of American party formation as an institutional solution to the above (asymmetric) coordination problem, whereby factions who were able to consolidate into larger cohorts enjoyed outsized political success compared to those who remained fragmented. This logic is furthered in the American historical lens by Bloch Rubin's \citeyear{rubin2013} examination of the ``House Insurgency'' between 1908 and 1910, in which initially-sporadic sparks of dissent were effectively organized into a coordinated movement largely (though not entirely) facilitated through party structures. Though suggestions of primarily-exogenous emergence provoked this further examination of endogenous formation, the argument had been well-established for quite some time; Lenin himself argued for the vanguard party as a means of overcoming coordination challenges which spontaneity had not been previously able to deliver \cite{lenin1935}.

It is further important to note that, largely because of this endogeneity, the construction of coordinating institutions does not necessarily take a uniform shape: even as these institutions reduce some friction between decision-makers, there can still be significant imbalances in the utilities achieved, largely resulting from agenda-setting powers. Romer \& Rosenthal's \citeyear{romerrosenthal1978} foundational examination of status quo reversion, for example, illustrates the privileges enjoyed by (1) the agenda-setter within a decision-making process, and (2) those who are favored under the status quo (as failure to collectively pass an alternative results in reversion to this policy preference). Their discussion is preceeded by the core social choice results of the 1970s now known as the McKelvey-Schofield Chaos Theorem(s) \cite{mckelvey1976,schofield1979}, which suggest that among multidimensional policy spaces, majority rule will not produce a unique Concorcet winner and almost any policy can be ultimately achieved through strategic agenda-setting. To this end, we potentially obtain a new situation in which further institutional constraints could provide enhanced stability -- but, we also can clearly see the perpetuation of noncooperative behavior even within institutionalized environments. Though it may be obvious, it is worth stating at this point -- institutionalization alone is not sufficient to solve cooperation. Rather, \emph{when} institutionalization is a possible solution, we must pursue the \emph{right} institutions.


\subsection*{Institutions for Collective Action}

A collective action problem exists when individual actors can be made better off by pursuing a collectively inferior outcome \cite{olson1965,taylor1990}. This may take the shape of resource overconsumption, such as the tragedy of the commons whereby every consumer of a collective resource has an incentive to overconsume \cite{olson1965,ostrom1990}; it may take the form of ``free-riding'' within an effort-needing task; and it can appear through a myriad of other manifestations, some of which we will explore later in this section and this response. At its core, however, the problem is simple: the pareto-efficient outcome is not an equilibrium, and therefore we cannot expect it to be supported by rational actors.

Discussing the foundations of institutions solving collective action problems is a bit more straightforward than with commitment problems, precisely because an institutional solution need only establish the efficient outcome as an equilibrium. Returning back to our construction of institutions as ``constraints on behavior,'' we can quickly recover two such applications. First, largely as a brute-force answer, institutions could potentially constrain action sets to only allow for the desired outcome, and in doing so fully eliminate the collective action problem. Second, and more realistically, institutions can be constructed to reduce incentives for non-cooperation by either increasing the value of total cooperation or providing some sanction against non-cooperative deviations. In either case, the role of the institution is, as Taylor \cite{taylor1990} describes, as an ``external'' solution: by changing the structure of the game, we can induce the desired behavioral outcome.

Yet, for a variety of reasons, it may not always be possible to eliminate collective-action defects or otherwise ``cheating'' behavior. In weak-institution environments, there may not be sufficient state capacity to establish such an external solution, or it may be the case that the supposed solution cannot be reasonably enforced: a fraud will not be deterred by the threat of government sanctions if the government is not able to reliably impose those sanctions. Here, institutions serve another role, and one that is particularly well-aligned with North's discussion of reducing transaction costs: institutions can limit the \emph{risk} of cooperative failures. Though we have already discussed that institutions may reduce market uncertainties, we can also conceptualize institutions which exist to mitigate the impact of those uncertainties, particularly when they are unavoidable. Enter, again, the law merchant \cite{milgrometal1990}! The expansion of trade throughout medieval times was burdened by the difficulty of enforcing exchange contracts in low-institutional environments and with traders from faraway markets. As Milgrom et al \citeyear{milgrometal1990} explain, a series of institutions both reduced the likelihood of a participant being cheated, and provided additional insurances were such a fraudulent interaction to occur. The result was a significant expansion in trade volumes, precisely because the imposition of locally-higher transaction costs (i.e., paying for the law merchant and enduring the institutional burden) allowed for globally-lower transaction costs and a lower-risk market. 

This follows a similar line of work on collective action institutions regarding the voluntary acquiescence of power and influence. I explored this rather substantially in my Core response, with respect to decisions by rulers and the elites to extend political franchises and constrain themselves from future exploits in an effort to nurture peace and investment. Though I will not re-hash that extensive line of reasoning as a whole, I again point to its relationship to reduced transaction costs by means of prohibiting collective action defection. As North \& Weingast \citeyear{northweingast1989} explore in seventeenth-century Britain, cooperative trade and investment followed from institutional restrictions on the ability of the crown to abuse property rights and renege on loans. In doing so, the state could effectively reduce the transaction costs experienced by private citizens who sought to accumulate and invest wealth, by meaningfully reducing the risks involved.


The emergence of institutions to solve collective action problems also appears relatively straightforward to justify: a group of actors experiencing collective action shortcomings (or otherwise finding their behavior hindered by the risk of such shortcomings) must ``simply" establish some institution(s) which either mitigates the incentives for defection or reduces the risk of harm from defection. When noncooperative conditions are particularly unpleasant, this can be a highly appealing process; indeed, Olson \citeyear{olson1965} in particular examines the establishment of labor unions as collective action institutions whose origin was (and still is) largely motivated by sufficiently poor working conditions and inequalities between labor and management. Yet, the prevalence of such institutions potentially obscures the existence of a second-order collective action problem, whereby the cost of \emph{creating} such institutions can induce free-riding in the hopes that others establish it instead \cite{hardin1990,ostrom1990}! It is precisely this conundrum that Sandler \citeyear{sandler1999} examines as a key reason for deviations from Olson's predictions, as the collective action problems arising from building institutions to solve collective action problems subsequently hinder their construction. It is potentially reasonable to suggest that creating third-order institutions to establish this institution-formation may be effective, and one could very well posit legislatures as holding such roles -- but what about the fourth-order, too? We simply face turtles all the way down\footnote{This was a running joke in Tom Clark's Winter 2025 Political Institutions course, when considering the endogenous construction of institutions and the endogenous formation of preferences. I feel obliged to my experience in that course to reference this accordingly.} in the hopes of some certain hearth of institutional formation, regardless of if one exists.



\subsection*{An Evaluatory Framework}

Though the above discussions are relatively foundational corners of social science, their inclusion in this response is not merely to demonstrate my knowledge of the literature\footnote{Though that's certainly a benefit!}. Rather, from these characterizations of the problems which institutions solve, we can construct a framework for evaluating whether or not institutional solutions are either necessary or optimal in establishing cooperative outcomes.

Let us again consider North's discussion of ``reduced transaction costs.'' If the role of institutions is to indeed facilitate lower-cost (and therefore more frequent) exchanges and interactions, then a core factor in whether or not a given institutional arrangement is necessary is whether or not it indeed provides lower transaction costs than would exist without it. North himself \citeyear{north1990a} examines this as a core evaluatory measure for institutional quality, positing that the effectiveness of an institution or arrangement of institutions can be broadly measured by evaluating its resultant impacts on transaction costs. Indeed, in considering the influence of the factors of institutional development explored in the main sections of this response, a core metric will be whether or not equivalent (or even lower) transaction costs can be achieved in their absence.

Yet, we must further examine exactly how we intend to consider transaction costs, for two reasons. First, as North \cite{north1990b} discusses, a subset of institutions may intentionally \emph{increase} some transaction costs. Various restrictions on property rights, for example, may be profitable for those with control over an institutional arrangement, and subsequently raise transaction costs.\footnote{A useful note here, however -- because this response is concerned with identifying when cooperation can be achieved with versus without institutions, the inclusion of ``malicious'' institutions can be almost summarily addressed: the alternative is almost certainly better.} Second, it may not always be clear exactly what ``altitude'' the salient transaction costs are being reduced; that is, we may often need to evaluate a much broader behavioral pathway in order to identify ultimate reduction, especially if the targeted reduction in cost is due to some externality. For example, many regimes place additional sales taxes on alcohol and tobacco products, which surely increases market friction and local costs between those markets' producers and consumers. Yet, in curbing alcohol and tobacco usage, the state may secure meaningful reductions in healthcare, emergency, and law enforcement expenses, freeing public funds for more productive investment. This \emph{broad} reduction in transaction cost is accordingly achieved through a \emph{localized} increase in transaction costs, but if we are to only consider the local lens, we may improperly evaluate these taxation institutions as ineffective. Similarly, we may find broad improvements in group cooperation achieved through local hindrances, and accordingly must be careful and intentional in considering the subsequent institutional impacts.

This points to a broader challenge in evaluating institutional design: because institutions are not generated randomly (more on that later), at least in their entirety, it can be difficult to determine their actual effects. Przeworski et al \citeyear{przeworski2000} provide an evaluatory framework based on counterfactuals that allow us to overcome this obstacle, at least in part. Their framework suggests that in order to properly analyze the impacts of institutions, we ``must ask what the outcomes would have been if other institutions had been observed under the same conditions'' \cite{przeworski2000}. This invites a degree of uncertainty in considering counterfactuals which we (tautologically) cannot observe, but is an important step forward in institutional evaluation nonetheless, and I make every effort to employ this framework in the remainder of this response.

I close this section, then, by focusing on an additional characterization of institutions: influences intended to steer interactions towards certain (ideally, socially-optimal) equilibria and away from others. An institution is an effective and necessary tool for cooperation, then, when two partially-overlapping conditions are met in equilibrium:

\begin{enumerate}
  \item Without the institution, a less-optimal equilibrium would be expected;
  \item No alternative (non-institutional) mechanism could support the present equilibrium
\end{enumerate}

Now that we have constructed the primary use cases of institutions within social, political, and economic life, we can use these conditions to evaluate the contexts under which they may not be essential -- or optimal -- for stability. The following sections explore the literature's primary ``exceptions'' whereby non-institutional alternatives may be best suited, and apply these conditions to generate suitable boundaries for necessary institutionalization.


\section*{Cooperation and Repetition}

Many collective action problems, including those discussed in the above sections, can be represented as a prisoners' dilemma. Such situations are based on a simple premise whereby joint cooperation is socially optimal, but under cooperation, each actor has an incentive to defect; subsequently, joint defection is socially sub-optimal. Though an extraordinarily common framework within the social sciences, an extensive game theoretic literature has focused on a non-institutional solution: repetition.

Suppose, in some collective action situation, that the actors involved will engage once and then depart, never to interact again. This is the implicit condition under which the one-shot prisoners' dilemma is conducted, and indeed represents the ``basic'' setting of collective action. Yet, Friedman \citeyear{friedman1971} demonstrated that a cooperative outcome can often be secured simply by repeating the game. In a conclusion now known as the ``Folk theorem'' or ``general feasibility theorem'' \cite{myerson1991}, Friedman proves that when actors repeatedly play a prisoners' dilemma and are sufficiently non-myopic, their ability to punish each other for defections can support an indefinitely-repeated cooperative equilibrium. The result provides for an environment where, in the absence of constraining institutions, the optimal equilibrium can in fact be achieved.

This framework is extraordinarily prevalent throughout the social sciences and is a core foundation of social cooperation. Take, for example, the alternation of power within democratic regimes. Elections are run by the state, meaning that the government has the \emph{capacity} for endogenous involvement within its own selection. With sufficient political support -- or simply sufficient vigor -- a ruling party can usually exert some degree of influence over electoral policies or processes in an effort to secure future political empowerment \cite{boix1999,diermeiervlaicu2011,krehbiel2004}. This degree of abuse can be particularly prevalent under majoritarian systems when a ruling party or actor can claim its ``mandate'' to reshape political life around the sentiments of their campaign \cite{powell2000}. Yet, given that effective democracy involves the legitimate alternation of power without abuse of electoral institutions \cite{przeworski2000}, some constraints on this ability seem appropriate. 

There are two such ways in which this abusive equilibrium (consolidation of power by whichever party is presently in office) can be avoided, and each provides insight into when institutional constraints may be most necessary. First, we might design electoral institutions such that they are incredibly difficult to change, and build institutions which indeed constraint the action set of the ruling party (at least as it pertains to future elections). This is a key construction of many liberal democracies, largely facilitated by means of constitutions and subsequent amendments: states may designate certain laws, such as those governing elections, as requiring enhanced thresholds for legislative alteration. This is indeed an effective way to curb abusive noncooperation, but it presents two challenges. First, as made clear by the expanse of the literature on institutional endogeneity, we want our institutions to be able to facilitate change \emph{when needed}. Throughout the history of modern democracy, such changes have often expanded the franchise, promoted political equality, and responded to technological advancements in voting and selection procedures. By enforcing significant institutional thresholds against such changes, we may prevent abuse, but similarly prevent non-abusive, productive developments. Second, though additional thresholds for electoral laws are quite common throughout modern democracies, they are not necessarily ubiquitous, at least in form. The United Kingdom, for example, has no single constitution, and though a broad range of resultant laws and judicial challenges would likely arise in the face of uni-party efforts to enact sweeping changes, there still exists a gap in institutionalization. Yet, the (relative) survival of democracy in the United Kingdom -- or at least the absence of any party dramatically rewriting its electoral structure in modern times -- is reflective of this reciprocative equilibrium. For one party to commit abuses would be to invite abuses against themselves when they next sat as the opposition.\footnote{The question is still open as to whether or not this equilibrium has been cast aside in the United States, but it is assuredly a relevant question. Efforts by the far-right to consolidate executive authority and reshape electoral and representative institutions in their own interest have become relatively commonplace since 2016, and in turn have invited some similar antidemocratic behaviors from the left. It is unclear whether or not this is enough to fully reflect a noncooperative equilibrium -- for example, a majority of mainstream Democratic legislators still seem hesitant to pursue broad institutional changes, including during the early Biden presidency when only the Supreme Court was majority-conservative -- yet, the propensity of such anti-institutional behavior has certainly increased nonetheless, likely indicating that historical aversions to pursuing such one-sided reforms were built on repetitive equilibria.}

As a whole, we can take the above to suggest that institutional solutions may be suitable for cooperation \emph{if they can construct repetition}. Under this logic, one solution to noncooperation may be establishing institutions which directly tie future possibilities to present actions. Again, the law merchant presents a suitable example, particularly alongside the period's system of private judges \cite{milgrometal1990}. The advent of long-distance trade meant that many merchants could suitably expect to meet each other only once, such that the future punishments of present exploitation may be relatively unimportant. While reputational mechanisms could have provided some reprieve -- a consistently dishonest trader will likely stir some broader conversation -- the new size of the trading communities inhibited the spread of consistent information between traders, at least at a speed fast enough to keep up with the missives. To alleviate this, Milgrom et al \citeyear{milgrometal1990} posits, the law merchant and affiliated private judge systems served not just as risk-reducing entities but also as informationally-binding institutions. The law merchant could keep a record of dishonest trades reported by his customers, such that a trade between two of his customers ensured\footnote{to a large extent} that (a) both had equal information, and (b) both were reputable enough to allow their behavior to be recorded. In effect, this created a ``buy-in'' community of repeated interactions whereby customers' present behavior would be intentionally presented in their future affairs, harnessing the incentive effects of repeated interactions despite most individual interactions remaining singular.

From this example, we can explore a further cornerstone of institutional emergence: namely (and unsurprisingly), institutions are likely to emerge when it is profitable for the institution's founder(s) to do so. For as much of an impact as the law merchant had on medeival trade, his presence was neither complimentary nor state-sanctioned. Rather, it became mutually productive to offer his services at a cost: his customers could still expect better returns from trade by minimizing their risk of being defrauded even less the law merchant's charge, and the law merchant would profit from providing his services. This reflects a broader discussion in institutional theory -- the notion of institutional stability stemming from self-enforcement. Under this line of logic, furthered by the likes of Ostrom \citeyear{ostrom1990} and Diermeier \& Krehbiel \citeyear{diermeierkrehbiel2003}, the enduring sustainability of an institutional arrangement is derived from the mutual benefit of the actors involved. When considering this as an equilibrium logic, the outcome is relatively clear -- if all parties are better off under the institutional arrangement than without it, none will have an incentive to pursue alternative arrangements \cite{northweingast1989}. This reasoning forms a crucial foundation of institutional emergence and perseverence, suggesting that choice-constraining institutions are most likely and most stable when those constraints ultimately benefit all actors involved.\footnote{There is a weaker application of this reasoning: perhaps only the salient actors need to be benefitted. It may be possible to provide institutions that do have ``losers,'' so long as those losers are not sufficiently empowered to either weaken or overthrow the institution itself. Though this is a less appealing condition -- complete mutual benefit is an entrancing aspect of institution-building -- it may also be significantly more realistic to achieve.}

However, Freidman's Folk theorem does not suggest the \emph{inevitability} of cooperation under repeated games, but rather, its possibility under sufficient discounting conditions \cite{friedman1971}. Namely, the theorem suggests that a cooperative equlibrium may be rationally maintained under repetition \emph{if and only if} the players sufficiently valued future payoffs. Highly myopic players, for example, may be relatively uninfluenced by payoffs from future interactions, such that the threat of future retribution from present noncooperation is not sufficient to secure the desired behavioral effects. This presents another consideration: institutions may indeed be necessary for cooperation if actors under repeated environments significantly discount future returns. To a large extent, we might expect to see this out of actors (including states) facing economic shortfalls or other immediate crises; if stable existence in the future is not secured without present defection, then consequences within that future existence are no longer meaningful constraints. This also can create difficulty amongst democratically-elected officials seeking to pursue cooperative outcomes: if there is a likelihood that either party will be replaced, and if that likelihood is in fact reduced by present non-cooperation, then we may expect to see shorter, weaker, or simply fewer cooperative terms. Here, repetition itself may not be sufficient to induce cooperation, and cooperation may instead require the institutional efforts discussed previously.

\section*{Pre-Existing Alignment and Informal Institutions}

Another line of work surrounds how social and contextual backgrounds may support cooperation in the absence of further institutionalization. However, this framework requires further delineating between two distinct arguments: that of \emph{exogenously existing} incentive alignment, and that of \emph{influential social norms}. The former is more concise; the latter is more interesting. 

In considering North's \citeyear{north1990} definition of institutions as ``the rules of the game,'' we can broadly characterize as an institution any such influence which shapes behavioral outcomes. Yet, this does not presume that the rules themselves are \emph{necessary} to deliver that outcome. In this sense, we can quickly return to our condition above, and state that institutions are not necessary for cooperation when that cooperation is (or would be) achieved and maintained without them. Take, for example, Diermeier \& Vlaicu's \citeyear{diermeiervlaicu2011} examination of party influence in the United States Congress. Under this structure, they find that parties are indeed effective institutions when beliefs and preferences are so far dispersed that collective legislating action could not be reliably achieved without such an enforcement mechanism; this falls in line with the discussion of parties above. Yet, they further find that when a group of legislators is sufficiently ideologically aligned, such that their potential gains from noncooperation on a coordinated action become rather minimal, they tend to achieve the cooperative outcome absent any party influence. This suggests that when existing incentives are sufficiently arranged and aligned, cooperation may indeed be achieved without furtner institutional involvement.

This conclusion is very closely tied to the above section's discussion of self-sustaining and self-enforcing institutions. It most closely reflects Shepsle's \& Boncheck \citeyear{shepslebonchek1997} characterization of collective cooperation in the absence of either coerceion or social expectations; along these lines, it is no surprise that we should then expect such cooperative endeavors to endure \emph{so long as the alignment of preferences does not significantly change}. This can potentially be a fatal flaw in form of cooperation, however, and makes it particularly vulnerable to paradigmatic changes. Just within the party context, for example, Harmel \& Janda \citeyear{harmeljanda1994} identify changes within common party goals as commonly resulting from either leadership change, a shift in factional dominance, or even just an external stimulus. If legislative cooperation without partisan organization is to be maintained purely by ideological alignment, then, we may very well want to keep the party's coercive capacity nearby and ready to employ.






\section*{Framing Norms and Reconciling Institutional Perspectives}


% =====================================================================
%  REFERENCES
%  Draws on two bibliographies at the repo root:
%    ../refs.bib    hand-curated. Edit freely; never overwritten.
%    ../zotero.bib  the whole Zotero library, generated by
%                   tools/export_zotero_bib.py. Do not edit by hand.
%  Look a key up in ../zotero-keys.txt, or just search ../zotero.bib.
%  Uncomment \nocite{*} to print every entry in refs.bib regardless of
%  whether it is cited — useful for checking the file, not for a draft.
% =====================================================================
\newpage
% \nocite{*}

% UNCOMMENT OUT ONCE STARTED
\bibliography{../refs,../zotero}

\end{document}
